\documentclass{article}
\usepackage[UTF8]{ctex}
\usepackage{amsmath,amssymb,amsfonts}
\usepackage{enumitem}
\usepackage{geometry}
\geometry{a4paper, margin=1in}
% === 错题本排版优化（不参与题目标识） ===
\setlength{\parindent}{0pt}
\setlength{\parskip}{0.35em plus 0.10em minus 0.08em}
\setlength{\emergencystretch}{3em}
\setlist{itemsep=0.25em, topsep=0.40em, parsep=0pt}
\allowdisplaybreaks[2]
% === 错题本排版优化结束 ===
\begin{document}
\section*{哈尔滨工业大学2024年高等代数真题}
\section*{74．哈尔滨工业大学2024年高等代数试卷}

\subsection{题目1}
1．解答如下问题：

\begin{itemize}
  \item[（1）] 证明：每个次数 \(\geq 3\) 的实系数多项式在实数域上一定可约。
  \item[（2）] 证明：三次实系数多项式在实数域上一定有根。
  \item[（3）] 四次实系数多项式在实数域上一定有根吗？说明理由。
\end{itemize}

\subsection{题目2}
2．讨论 \(a, b\) 取何值时下列方程组有解，并求解。

\[
\left\{\begin{array}{l}
x_{1}+x_{2}+x_{3}+x_{4}+x_{5}=1 \\
3 x_{1}+2 x_{2}+x_{3}+x_{4}+a x_{5}=-3 \\
x_{2}+2 x_{3}+2 x_{4}+6 x_{5}=3 \\
5 x_{1}+4 x_{2}+3 x_{3}+3 x_{4}-x_{5}=b
\end{array}\right.
\]

\subsection{题目3}
3．设 \(\eta\) 是非齐次线性方程组 \(A X=\beta(\beta \neq 0)\) 的一个解向量，\(\xi_{1}, \xi_{2}, \cdots, \xi_{t}\) 是对应齐次线性方程组\(A X=0\) 的一个基础解系，证明：

\begin{itemize}
  \item[（1）] \(\eta, \xi_{1}, \xi_{2}, \cdots, \xi_{t}\) 线性无关。
  \item[（2）] \(\eta, \xi_{1}+\eta, \xi_{2}+\eta, \cdots, \xi_{t}+\eta\) 是 \(A X=\beta\) 的线性无关的解向量．
  \item[（3）] \(A X=\beta\) 的任意解 \(Y\) 都可以表示成
\[
Y=k_{0} \eta+k_{1}\left(\xi_{1}+\eta\right)+k_{2}\left(\xi_{2}+\eta\right)+\cdots+k_{t}\left(\xi_{t}+\eta\right) .
\]
其中 \(k_{0}+k_{1}+\cdots+k_{t}=1\) 。

\subsection{题目4}
\end{itemize}

\subsection{题目4}
原题边界“4．”：设 \(V=\left\{\left(x_{1}, x_{2}, x_{3}, x_{6}, x_{5}\right)^{\mathrm{T}} \in \mathbb{R}^{6} \mid x_{1}+7 x_{2}+5 x_{3}-4 x_{4}+2 x_{5}=0\right\}\) 。

\begin{itemize}
\item[（1）] 证明：\(S=\left\{(-2,0,0,-1,-1)^{\mathrm{T}},(1,1,-2,-1,-1)^{\mathrm{T}},(-5,1,0,1,1)^{\mathrm{T}}\right\}\) 是 \(V\) 的一个线性无关的子集．
  \item[（2）] 将 \(S\) 扩充为 \(V\) 的一组基底。
\end{itemize}

\subsection{题目5}
5．设 \(\alpha_{1}, \alpha_{2}, \cdots, \alpha_{r}\) 是线性无关的 \(n\) 维列向量，\(\beta_{i}=\sum_{j=1}^{r} a_{i j} \alpha_{j}, i=1,2, \cdots, r\) ，证明：\(\beta_{1}, \beta_{2}, \cdots, \beta_{r}\) 线性相关的充要条件为

\[
\left|\begin{array}{cccc}
a_{11} & a_{12} & \cdots & a_{1 r} \\
a_{21} & a_{22} & \cdots & a_{2 r} \\
\vdots & \vdots & & \vdots \\
a_{r 1} & a_{r 2} & \cdots & a_{r r}
\end{array}\right|=0 .
\]

\subsection{题目6}
6．设 \(A\) 为 \(n\) 阶方阵，\(\lambda_{0}\) 是 \(A\) 的几何重数为 \(m\) 的特征值，证明：\(\left(\lambda-\lambda_{0}\right)^{m} \| \lambda E_{n}-A \mid\) 。

\subsection{题目7}
7．设 \(P^{n}\) 是数域 \(P\) 上的 \(n\) 维线性空间，\(\alpha_{1}, \alpha_{2}, \cdots, \alpha_{s}\) 和 \(\beta_{1}, \beta_{2}, \cdots, \beta_{t}\) 是 \(P^{n}\) 中的两组向量。

\begin{itemize}
  \item[（1）] 给出并证明 \(L\left(\alpha_{1}, \alpha_{2}, \cdots, \alpha_{n}\right) \cup L\left(\beta_{1}, \beta_{2}, \cdots, \beta_{t}\right)\) 是 \(P^{n}\) 的子空间的充要条件。
  \item[（2）] 给出并证明 \(L\left(\alpha_{1}, \alpha_{2}, \cdots, \alpha_{n}\right) \cup L\left(\beta_{1}, \beta_{2}, \cdots, \beta_{t}\right)=P^{n}\) 的充要条件。
\end{itemize}

\subsection{题目8}
8．设 \(A\) 是实对称矩阵，\(X=\left(x_{1}, x_{2}, x_{3}, x_{4}\right)^{\mathrm{T}} \in \mathbb{R}^{4}, f(X)=X^{\mathrm{T}} A X\) 是四元实二次型．

\begin{itemize}
  \item[（1）] \(S=\left\{X \in \mathbb{R}^{4} \mid f(X)=0\right\}\) 是不是 \(\mathbb{R}^{4}\) 的子空间？为什么？
  \item[（2）] 若矩阵 \(A\) 的正负惯性指数都是 1 ，证明：存在 \(\mathbb{R}^{3}\) 的 3 维子空间 \(W\) ，使得当 \(X \in W\) 时，\(f(X)=0\)
\end{itemize}

\subsection{题目9}
9．设 \(V\) 是 \(n\) 维复线性空间， \(\mathcal{A}, \mathcal{B}\) 是 \(V\) 的线性变换， \(\mathcal{A} \mathcal{B}=\mathcal{B} \mathcal{A}, f(t)\) 是一个复系数多项式。证明：

\begin{itemize}
  \item[（1）] \(\mathcal{B}\) 的特征子空间与核空间的和是 \(\mathcal{A}\) 的不变子空间。
  \item[（2）] \(f(\mathcal{A})\) 和 \(f(\mathcal{B})\) 有共同的特征向量．
\end{itemize}

\subsection{题目10}
10．已知

\[
A=\left(\begin{array}{lllll}
0 & 0 & 0 & 0 & 0 \\
1 & 0 & 0 & 0 & 1 \\
0 & 1 & 0 & 0 & 0 \\
0 & 0 & 1 & 0 & 0 \\
0 & 0 & 0 & 1 & 0
\end{array}\right)
\]

求 \(A\) 的不变因子和 Jordan 标准形．
\end{document}
